% GENERATED FILE. Edit canonical lessons, then run npm run generate:books.
% canonical-source-hash: fnv1a64:34248da51a0a692e
% canonical-lessons: TA-C505-ayul, TA-C505-maranam, TA-C505-uyil, TA-C505-vitavai, TA-C505-anatai, TA-R505-review-499-505, TA-R505-review-final-505

\chapter{The End of Life}
\label{ch:ta-the-end-of-life}
\hlchaptermodality{505}

\begin{chapteropening}
\textbf{By the end of this chapter:} \emph{I can talk about a lifespan, death, wills, widows and orphans.}

Complete the last lesson of chapter 505: I can talk about a lifespan, death, wills, widows and orphans.
\end{chapteropening}

\section[\texorpdfstring{\ta{ஆயுள்} (āyuḷ)}{ஆயுள் (āyuḷ)}]{\ta{ஆயுள்} (āyuḷ) — a lifespan}

\label{lesson:TA-C505-ayul}



\begin{quote}
Before the new one: say the Tamil for to fade, then the Tamil for to shrink.
\end{quote}

\subsection*{You'll want to know: \ta{ஆயுள்}}
\textbf{\ta{ஆயுள்}} — \emph{āyuḷ} — \textquotedblleft{}a lifespan\textquotedblright{}.

\begin{quote}
\ta{நல்ல} \ta{உணவு} \ta{ஆயுளை} \ta{அதிகரிக்கும்}.
\end{quote}

\begin{quote}
\emph{nalla uṇavu āyuḷai atikarikkum.}
\end{quote}

\begin{quote}
Good food lengthens one’s life.
\end{quote}

\ta{ஆயுள்} \ta{முழுவதும்}, all one’s life.

\begin{cousinweb}
\textbf{\ta{ஆயுள்}} (\emph{āyuḷ}) is usually traced to Sanskrit \emph{āyus}, \textquotedblleft{}life, the length of life\textquotedblright{}; our sources do not list the Tamil form.
\end{cousinweb}

\begin{grammarlens}[title={What you've built}]
One more word for describing what you lived and felt.
\end{grammarlens}

\subsection*{Your turn}
\begin{itemize}
\raggedright
  \item \emph{Say it:} \emph{āyuḷ}
  \item \emph{Say it:} \emph{āyuḷ}, once more
  \item \emph{Say it:} \emph{nalla uṇavu āyuḷai atikarikkum.} — Good food lengthens one’s life.
  \item \emph{Recall:} say the Tamil for to fade, then the Tamil for to shrink, then say \emph{āyuḷ} again
\end{itemize}

\begin{tcolorbox}[breakable,colback=teal!4,colframe=teal!35!black,title={Before you move on}]
What does \ta{ஆயுள்} mean? (\textquotedblleft{}A lifespan\textquotedblright{}.) Say it once more.
\end{tcolorbox}
\section[\texorpdfstring{\ta{மரணம்} (maraṇam)}{மரணம் (maraṇam)}]{\ta{மரணம்} (maraṇam) — death}

\label{lesson:TA-C505-maranam}



\begin{quote}
Before the new one: say the Tamil for to shrink, then the Tamil for a lifespan.
\end{quote}

\subsection*{You'll want to know: \ta{மரணம்}}
\textbf{\ta{மரணம்}} — \emph{maraṇam} — \textquotedblleft{}death\textquotedblright{}.

\begin{quote}
\ta{அவருடைய} \ta{மரணம்} \ta{எல்லோருக்கும்} \ta{அதிர்ச்சியாக} \ta{இருந்தது}.
\end{quote}

\begin{quote}
\emph{avaruṭaiya maraṇam ellōrukkum atircciyāka iruntatu.}
\end{quote}

\begin{quote}
His death was a shock to everyone.
\end{quote}

Formal and written; in speech \ta{சாவு}.

\begin{cousinweb}
\textbf{\ta{மரணம்}} (\emph{maraṇam}) is borrowed from Sanskrit \emph{maraṇa}, \textquotedblleft{}dying\textquotedblright{} (CDIAL 9869).
\end{cousinweb}

\begin{grammarlens}[title={What you've built}]
One more word for describing what you lived and felt.
\end{grammarlens}

\subsection*{Your turn}
\begin{itemize}
\raggedright
  \item \emph{Say it:} \emph{maraṇam}
  \item \emph{Say it:} \emph{maraṇam}, once more
  \item \emph{Say it:} \emph{avaruṭaiya maraṇam ellōrukkum atircciyāka iruntatu.} — His death was a shock to everyone.
  \item \emph{Recall:} say the Tamil for to shrink, then the Tamil for a lifespan, then say \emph{maraṇam} again
\end{itemize}

\begin{tcolorbox}[breakable,colback=teal!4,colframe=teal!35!black,title={Before you move on}]
What does \ta{மரணம்} mean? (\textquotedblleft{}Death\textquotedblright{}.) Say it once more.
\end{tcolorbox}
\section[\texorpdfstring{\ta{உயில்} (uyil)}{உயில் (uyil)}]{\ta{உயில்} (uyil) — a will}

\label{lesson:TA-C505-uyil}



\begin{quote}
Before the new one: say the Tamil for a lifespan, then the Tamil for death.
\end{quote}

\subsection*{You'll want to know: \ta{உயில்}}
\textbf{\ta{உயில்}} — \emph{uyil} — \textquotedblleft{}a will\textquotedblright{}.

\begin{quote}
\ta{தாத்தா} \ta{தன்} \ta{வீட்டை} \ta{உயிலில்} \ta{பேரனுக்கு} \ta{எழுதி} \ta{வைத்தார்}.
\end{quote}

\begin{quote}
\emph{tāttā taṉ vīṭṭai uyilil pēraṉukku eḻuti vaittār.}
\end{quote}

\begin{quote}
Grandfather left his house to his grandson in his will.
\end{quote}

\ta{உயில்} \ta{எழுது}, to make a will.

\begin{cousinweb}
\textbf{\ta{உயில்}} (\emph{uyil}) is a loan from English \emph{will}, reshaped to Tamil sounds.
\end{cousinweb}

\begin{grammarlens}[title={What you've built}]
One more word for describing what you lived and felt.
\end{grammarlens}

\subsection*{Your turn}
\begin{itemize}
\raggedright
  \item \emph{Say it:} \emph{uyil}
  \item \emph{Say it:} \emph{uyil}, once more
  \item \emph{Say it:} \emph{tāttā taṉ vīṭṭai uyilil pēraṉukku eḻuti vaittār.} — Grandfather left his house to his grandson in his will.
  \item \emph{Recall:} say the Tamil for a lifespan, then the Tamil for death, then say \emph{uyil} again
\end{itemize}

\begin{tcolorbox}[breakable,colback=teal!4,colframe=teal!35!black,title={Before you move on}]
What does \ta{உயில்} mean? (\textquotedblleft{}A will\textquotedblright{}.) Say it once more.
\end{tcolorbox}
\section[\texorpdfstring{\ta{விதவை} (vitavai)}{விதவை (vitavai)}]{\ta{விதவை} (vitavai) — a widow}

\label{lesson:TA-C505-vitavai}



\begin{quote}
Before the new one: say the Tamil for death, then the Tamil for a will.
\end{quote}

\subsection*{You'll want to know: \ta{விதவை}}
\textbf{\ta{விதவை}} — \emph{vitavai} — \textquotedblleft{}a widow\textquotedblright{}.

\begin{quote}
\ta{கணவர்} \ta{இறந்த} \ta{பிறகு} \ta{அவள்} \ta{விதவை} \ta{ஆனாள்}.
\end{quote}

\begin{quote}
\emph{kaṇavar iṟanta piṟaku avaḷ vitavai āṉāḷ.}
\end{quote}

\begin{quote}
After her husband died she became a widow.
\end{quote}

Formal; \ta{கைம்பெண்} is the native word.

\begin{cousinweb}
\textbf{\ta{விதவை}} (\emph{vitavai}) is borrowed from Sanskrit \emph{vidhavā}, \textquotedblleft{}widow\textquotedblright{} (CDIAL 11752); the final -ā becomes Tamil -\ta{ஐ}.
\end{cousinweb}

\begin{grammarlens}[title={What you've built}]
One more word for describing what you lived and felt.
\end{grammarlens}

\subsection*{Your turn}
\begin{itemize}
\raggedright
  \item \emph{Say it:} \emph{vitavai}
  \item \emph{Say it:} \emph{vitavai}, once more
  \item \emph{Say it:} \emph{kaṇavar iṟanta piṟaku avaḷ vitavai āṉāḷ.} — After her husband died she became a widow.
  \item \emph{Recall:} say the Tamil for death, then the Tamil for a will, then say \emph{vitavai} again
\end{itemize}

\begin{tcolorbox}[breakable,colback=teal!4,colframe=teal!35!black,title={Before you move on}]
What does \ta{விதவை} mean? (\textquotedblleft{}A widow\textquotedblright{}.) Say it once more.
\end{tcolorbox}
\section[\texorpdfstring{\ta{அனாதை} (aṉātai)}{அனாதை (aṉātai)}]{\ta{அனாதை} (aṉātai) — an orphan}

\label{lesson:TA-C505-anatai}



\begin{quote}
Before the new one: say the Tamil for a will, then the Tamil for a widow.
\end{quote}

\subsection*{You'll want to know: \ta{அனாதை}}
\textbf{\ta{அனாதை}} — \emph{aṉātai} — \textquotedblleft{}an orphan\textquotedblright{}.

\begin{quote}
\ta{அனாதைக்} \ta{குழந்தைகளுக்கு} \ta{அவர்} \ta{உதவினார்}.
\end{quote}

\begin{quote}
\emph{aṉātaik kuḻantaikaḷukku avar utaviṉār.}
\end{quote}

\begin{quote}
He helped orphaned children.
\end{quote}

\ta{அனாதை} \ta{இல்லம்}, an orphanage.

\begin{cousinweb}
\textbf{\ta{அனாதை}} (\emph{aṉātai}) is borrowed from Sanskrit \emph{anātha}, \textquotedblleft{}having no master or protector\textquotedblright{} (Monier-Williams).
\end{cousinweb}

\begin{grammarlens}[title={What you've built}]
One more word for describing what you lived and felt.
\end{grammarlens}

\subsection*{Your turn}
\begin{itemize}
\raggedright
  \item \emph{Say it:} \emph{aṉātai}
  \item \emph{Say it:} \emph{aṉātai}, once more
  \item \emph{Say it:} \emph{aṉātaik kuḻantaikaḷukku avar utaviṉār.} — He helped orphaned children.
  \item \emph{Recall:} say the Tamil for a will, then the Tamil for a widow, then say \emph{aṉātai} again
\end{itemize}

\begin{tcolorbox}[breakable,colback=teal!4,colframe=teal!35!black,title={Before you move on}]
What does \ta{அனாதை} mean? (\textquotedblleft{}An orphan\textquotedblright{}.) Say it once more.
\end{tcolorbox}
\section[(dialogue)]{Looking back at family and the stages of life}

\label{lesson:TA-R505-review-499-505}



\begin{quote}
Say the Tamil for a dowry and a blessing.
\end{quote}

\subsection*{Your turn}
Say these aloud:

\begin{itemize}
\raggedright
  \item a dowry, a blessing, to marry, a married couple, married life
  \item a companion, a homemaker, a woman’s marital home, a separate household, to get along
  \item a squabble, a dispute, ill feeling, to sulk, a divorce
  \item to accommodate, to come close, to embrace, to stroke, a bond
  \item a female friend, a stranger, an introduction, to host, a farewell
  \item old age, an elderly person, memory, to fade, to shrink
  \item a lifespan, death, a will, a widow, an orphan
\end{itemize}

\begin{tcolorbox}[breakable,colback=teal!4,colframe=teal!35!black,title={Before you move on}]
A dowry? (\textbf{\ta{வரதட்சணை}}, \emph{varataṭcaṇai}.) A blessing? (\textbf{\ta{ஆசீர்வாதம்}}, \emph{ācīrvātam}.) To marry? (\textbf{\ta{மண}}, \emph{maṇa}.) A married couple? (\textbf{\ta{தம்பதி}}, \emph{tampati}.) Married life? (\textbf{\ta{இல்லறம்}}, \emph{illaṟam}.) A companion? (\textbf{\ta{துணை}}, \emph{tuṇai}.) A homemaker? (\textbf{\ta{இல்லத்தரசி}}, \emph{illattaraci}.) A woman’s marital home? (\textbf{\ta{புகுந்த} \ta{வீடு}}, \emph{pukunta vīṭu}.) A separate household? (\textbf{\ta{தனிக்குடித்தனம்}}, \emph{taṉikkuṭittaṉam}.) To get along? (\textbf{\ta{ஒத்துப்போ}}, \emph{ottuppō}.) A squabble? (\textbf{\ta{சச்சரவு}}, \emph{caccaravu}.) A dispute? (\textbf{\ta{தகராறு}}, \emph{takarāṟu}.) Ill feeling? (\textbf{\ta{மனஸ்தாபம்}}, \emph{maṉastāpam}.) To sulk? (\textbf{\ta{பிணங்கு}}, \emph{piṇaṅku}.) A divorce? (\textbf{\ta{விவாகரத்து}}, \emph{vivākarattu}.) To accommodate? (\textbf{\ta{அனுசரி}}, \emph{aṉucari}.) To come close? (\textbf{\ta{நெருங்கு}}, \emph{neruṅku}.) To embrace? (\textbf{\ta{தழுவு}}, \emph{taḻuvu}.) To stroke? (\textbf{\ta{வருடு}}, \emph{varuṭu}.) A bond? (\textbf{\ta{பந்தம்}}, \emph{pantam}.) A female friend? (\textbf{\ta{தோழி}}, \emph{tōḻi}.) A stranger? (\textbf{\ta{அந்நியர்}}, \emph{anniyar}.) An introduction? (\textbf{\ta{அறிமுகம்}}, \emph{aṟimukam}.) To host? (\textbf{\ta{உபசரி}}, \emph{upacari}.) A farewell? (\textbf{\ta{பிரியாவிடை}}, \emph{piriyāviṭai}.) Old age? (\textbf{\ta{முதுமை}}, \emph{mutumai}.) An elderly person? (\textbf{\ta{முதியவர்}}, \emph{mutiyavar}.) Memory? (\textbf{\ta{ஞாபகம்}}, \emph{ñāpakam}.) To fade? (\textbf{\ta{மங்கு}}, \emph{maṅku}.) To shrink? (\textbf{\ta{சுருங்கு}}, \emph{curuṅku}.) A lifespan? (\textbf{\ta{ஆயுள்}}, \emph{āyuḷ}.) Death? (\textbf{\ta{மரணம்}}, \emph{maraṇam}.) A will? (\textbf{\ta{உயில்}}, \emph{uyil}.) A widow? (\textbf{\ta{விதவை}}, \emph{vitavai}.) An orphan? (\textbf{\ta{அனாதை}}, \emph{aṉātai}.)
\end{tcolorbox}
\section[(dialogue)]{The End of Life — once more}

\label{lesson:TA-R505-review-final-505}



\begin{quote}
Say the Tamil for a lifespan and death.
\end{quote}

\subsection*{Your turn}
\begin{itemize}
\raggedright
  \item \emph{Say it:} a lifespan, death, a will, a widow, an orphan
\end{itemize}

\begin{tcolorbox}[breakable,colback=teal!4,colframe=teal!35!black,title={Before you move on}]
A lifespan? (\textbf{\ta{ஆயுள்}}, \emph{āyuḷ}.) Death? (\textbf{\ta{மரணம்}}, \emph{maraṇam}.) A will? (\textbf{\ta{உயில்}}, \emph{uyil}.) A widow? (\textbf{\ta{விதவை}}, \emph{vitavai}.) An orphan? (\textbf{\ta{அனாதை}}, \emph{aṉātai}.)
\end{tcolorbox}
